\section{Interpretation and limitations}
\label{sec:interpretation}

\subsection{What the exercise establishes}

The narrowest statement of the result is also the firmest. A valuation model built
from public filings and public curve data, calibrated entirely on 31~December 2021
information and then frozen, priced two large securities portfolios a year forward
with errors of two to four per cent, and reproduced the economic-value profile SVB
published for 2021 at shock points it was not fitted to, including a feature of that
profile's shape. Nothing in the construction required data unavailable at the
calibration date, a view on the path of rates, or hindsight about either bank.

From that follows the measurement claim the paper is for. On the deposit assumption
implied by SVB's own published sensitivity, SVB's worst-case loss under the
standardised scenarios at the end of 2021 was 36.2\% of Tier~1 capital, against a
supervisory threshold of 15\%. That figure is derivable from the FY2021 filing
alone, fifteen months before the bank failed. The exposure was measurable from what
had already been disclosed.

\subsection{What the deposit assumption actually was}

The more consequential reading concerns what a disclosed economic-value figure is.

The asymmetry set out in Section~\ref{sec:sensitivity} has a consequence for how such
a figure should be read. Securities are reported at fair value by type and maturity,
and Section~\ref{sec:comparison} shows that side of the calculation can be rebuilt
and checked against subsequent disclosure. The deposit side cannot be checked against
anything, because neither bank stated the assumption behind it. That
assumption moves the measured exposure across a range wider than the exposure itself.

So SVB's disclosed $-13.9\%$ at $+100$ basis points was not a measurement in the way
its disclosed securities fair values were. It was the output of a model containing at
least one assumption that the filing did not state and that no reader could have
recovered. Solving for the value consistent with the disclosure gives an uninsured
half-life of 3.59~years. For a deposit base that was 96\% uninsured and concentrated
in venture-backed technology firms, that is a claim about depositor loyalty rather
than a property of the balance sheet.

This interpretation is the paper's, not a finding. What the evidence supports is
narrower and should be kept distinct: the disclosed sensitivity is reproducible under
a 3.59-year half-life (established), the measured exposure is highly sensitive to that
parameter (established), and nothing in the filings pins it down (established). That
the assumption was implausible for this particular depositor base is a judgement,
supported by the composition figures in Section~\ref{sec:objects} and by the mechanism
in \citet{Drechsler2021}, under which a deposit franchise hedges a long asset book
only for as long as the deposits stay. \citet{Jiang2024} identify uninsured leverage
as the variable that determines whether mark-to-market losses become solvency
problems, which is the same point from the other direction.

\subsection{A conflation in the implied half-life}

One qualification belongs here rather than in a list, because it affects how the
3.59-year figure should be read.

SVB's FY2021 filing states that the economic-value exposures in its sensitivity table
include the effect of its hedging, and that the $+100$ basis point figure would have
been $-15.4\%$ without any hedging applied, reduced to $-13.9\%$ by pay-fixed swaps
purchased in March 2021 \citep{SVB10K2021}. The reconstruction in this paper does not
model that swap book.

The half-life solved in Section~\ref{sec:comparison} is therefore the deposit
assumption that reproduces a \emph{post-hedge} sensitivity using an \emph{unhedged}
balance sheet, and it absorbs whatever the swaps were contributing.

The size of that contamination can be measured rather than conceded, because SVB
discloses the unhedged figure in the same filing. Re-solving against $-15.4\%$ puts
an unhedged model against an unhedged disclosure, and returns an uninsured half-life
of 3.42~years against the 3.59 obtained from the published figure. The conflation is
worth 0.17~years. Set against a parameter whose plausible range in
Section~\ref{sec:sensitivity} spans six months to ten years, and whose movement from
3.59 to 2.00 years more than doubles the measured exposure, 0.17~years changes
nothing that matters.

A second result from the same solve is more interesting, because it concerns a
quantity that was never fitted. Calibrated post-hedge, the model returns a base
economic value of \$21,271m against SVB's disclosed \$20,672m, an overstatement of
2.90\%. Calibrated like-for-like against the pre-hedge figure, it returns \$20,808m,
an overstatement of 0.66\%. Adjusting the calibration to remove a known mismatch
improved the accuracy of a level the calibration does not target, by a factor of four.
That is weak evidence, since a single quantity moving the right way proves little, but
it is evidence in the direction the specification predicts.

What the pre-hedge solve does not provide is a replacement validation. Its $+200$
basis point figure is $-30.40\%$, which cannot be set against SVB's disclosed
$-27.70\%$ because the disclosure is post-hedge and this figure is not. The published
profile and the published pre-hedge number are on different bases, and only one of
them extends beyond $+100$ basis points. The paper therefore keeps the post-hedge
calibration for the out-of-sample comparison in Section~\ref{sec:comparison}, where
both sides are on the same basis, and reports the pre-hedge solve here as the cleaner
estimate of the deposit parameter itself. Modelling the swap book directly would
collapse the two, and the filing discloses the notional (\$10,700m) and fair value
(\$18m) but neither the tenor nor the fixed rate, so it is not something this
reconstruction can do without inventing an assumption of exactly the kind the paper
objects to.

The hedging itself is worth one further note. \citet{Kim2024} finds that banks did not
systematically ignore balance-sheet risk in this period, increasing hedging as
securities losses grew and scaling it back as they reversed, and concludes that SVB's
conduct was idiosyncratic rather than representative. SVB did hold pay-fixed swaps,
and its own filing records them reducing the $+100$ basis point exposure by 1.5
percentage points. The reconstruction here is silent on whether that was adequate,
and nothing in this paper should be read as evidence on the question.

\subsection{Why one bank failed and the other did not}

The two banks reached the end of 2022 holding unrecognised losses within a billion
dollars of each other, as Table~\ref{tab:objects} records, and both breach the
supervisory outlier screen in every bank-year tested. On the run model, only SVB
crosses a threshold, and it crosses on liquidity rather than solvency; Schwab crosses
neither in either year under any of the 42 parameter combinations tried.

That pattern is consistent with the uninsured-leverage mechanism of
\citet{Jiang2024}, and the paper reports it as consistency rather than as
identification. Two banks differing in uninsured share also differ in depositor type,
business model, asset mix, size and supervisory treatment, all at once. A design with
two observations cannot attribute the difference in run thresholds to any one of
them, and the sentence that the funding structure \emph{is} what separates these
banks would require a cross-section this paper does not have. What can be said is
narrower: holding the modelling framework fixed across both, the threshold moves with
the uninsured share and not with the size of the securities loss.

The limits run further than sample size. \citet{KellyRose2025} argue that the 2023
failures are better explained by the business models of the banks involved, in SVB's
case its concentration in venture-capital-backed clients whose own funding had
deteriorated, than by balance-sheet metrics, and that this account separates the
affected banks more precisely. Nothing in the present design adjudicates that. A
reconstruction establishes what the filings could reveal about exposure; it does not
establish which variable best predicts failure, and the two questions have been
treated here as distinct throughout. Whether funding structure separated the actual
outcomes in March 2023 is in any case beyond this design: SVB failed amid a
concentrated, socially coordinated withdrawal that no static balance-sheet model
represents, and the counterfactual in which it had Schwab's funding mix is not
observable.

\subsection{What the paper does not establish}

Three claims are specifically not made.

It is not claimed that the absence of an economic-value figure from SVB's FY2022
filing contributed to the failure. The paper establishes that the exposure was
measurable from the FY2021 filing; it does not test what any party would have done
with a 2022 figure, and the supervisory record indicates that the direction of the
problem was already recognised \citep{FRB2023}.

It is not claimed that SVB's disclosed figures were incorrect or misleading. The
disclosed profile is internally consistent and this paper reproduces it. The argument
concerns what such a figure can convey when the assumption driving it is not stated,
not whether this one was computed honestly.

And no statistical claim is made anywhere. The results are deterministic re-solves of
a model; the ranges reported are ranges across assumptions, not sampling
distributions. Nothing in this paper should be read as a significance test or a
confidence interval.

\subsection{Limitations}

The ones that would change a reported number, rather than a list of everything
imperfect.

\textit{Book yields stand in for coupons.} A bond bought at a premium is approximated
by a par bond of the same yield. The option-adjusted spread absorbs the level error;
the residual duration difference is not corrected.

\textit{The prepayment curve is stylised.} Its parameters are not estimated from
pool-level data, because the filings disclose neither loan-level balances nor actual
weighted-average coupons. They are held fixed across both banks and both dates so
nothing can be tuned after the answer is visible, and
Section~\ref{sec:sensitivity} shows them to be the largest single source of
residual error.

\textit{Representative maturities stand in for contractual buckets.} The open-ended
over-ten-years bucket is modelled at twenty years for bullets and twenty-five for
pass-throughs. Varying this moves the headline prediction by about a point in either
direction.

\textit{Mean reversion is weakly identified.} The best $R^2$ in any specification
tried is 0.0387. This is tolerable only because the reported quantity barely depends
on the parameter, which Section~\ref{sec:sensitivity} demonstrates rather than
asserts.

\textit{One bucket failed to calibrate.} SVB's \$100m of variable-rate agency CMOs is
disclosed above par and the solver cannot reach that level for a floating-rate
position. The shortfall is \$0.27m on a \$126bn book.

\textit{Loans carry no credit spread.} The filings disclose no loan-level fair value
to calibrate one against. This understates credit risk and does not affect the
rate sensitivity the economic-value measure captures.

\textit{Schwab's yields are bucket-level.} Schwab reports yields by maturity band
rather than by security type, so its positions inherit a band-level yield where SVB's
inherit a type-level one.

\textit{The run model has no clock.} It is a static balance-sheet calculation. It
says what the balance sheet could absorb, not how fast a run moves, which is why the
67.5\% threshold is better read as an upper bound than as a forecast.

\textit{The 2024 shock recalibration is reported but not adopted.} It takes effect
from January 2026 and post-dates both valuation dates. Its effect is reported in
Section~\ref{sec:sensitivity}, where it changes which scenario binds for SVB without
changing any outlier-test verdict.
